\section{Verification, and what a verdict is worth}
\label{sec:verification}

\subsection{The separation of concerns that took longest to get right}

For most of this project's life, ``verify a solve'' meant ``reconstruct the
move sequence exactly enough that it provably reaches solved''. That framing
is wrong, and the reason is worth stating carefully because it is the
conceptual core of the whole system.

Reconstruction accuracy and legitimacy are different questions, and the
first is a \emph{terrible} instrument for the second.

\begin{keybox}
\textbf{Verification is a cliff.} Empirically, sessions above ${\sim}98\%$
raw per-move accuracy reconstruct exactly and verify; sessions in the 95--98\%
band verify essentially never (0 of 4 in the recorded corpus). There is no
graded middle. So a legitimacy test built on exact reconstruction rejects
honest users for being filmed in slightly worse light --- it is a
\emph{model-quality} test wearing a legitimacy test's clothes.
\end{keybox}

The system therefore splits into two tiers with different jobs:

\begin{description}[leftmargin=1.2em,style=nextline,itemsep=3pt]
\item[The accuracy tier (the group decode).] Makes the recorded transcript as
close to what was performed as possible. Its output is a product feature
(coaching, move-by-move review). It issues no verdicts.
\item[The legitimacy tier (the count gate).] Decides
\texttt{VERIFIED}/\texttt{REVIEW}/\texttt{REJECTED}. It uses only the
\emph{number} of moves, never their identity.
\end{description}

\subsection{Falsifiability: the measurement that makes a verdict mean
anything}
\label{sec:falsifiability-method}

A verifier that has only ever been run on genuine solves has not been
evaluated. \texttt{VERIFIED} is not a result until you know what the same
video says to a claim that is false.

So every decode is re-run against deliberately wrong claims, under an
\emph{identical} cost model: the same onsets, the same posteriors, the same
beam, decoded against a start state that is 1, 2 or 4 quarter turns away
from the truth, and against a cube that was never scrambled at all. A
verifier that accepts these is worthless; one that rejects them has
demonstrated something.

\paragraph{And here is the part that must not be oversold.} Suppose the true
claim decodes to word $W$ at cost $c(W)$. A claim whose start state is $d$
quarter turns away from the truth is solved by the word $w^{-1}W$, where $w$
is the offending $d$-move prefix. That word accepts exactly the same onsets
as $W$ and additionally performs $d$ insertions. So a valid story for the
false claim \textbf{provably exists} at cost
\begin{equation}
c(W) + d\cdot C_{\mathrm{ins}},
\label{eq:decoybound}
\end{equation}
whether or not the beam finds it. When the system rejects a 1-move-off
decoy, it is being rejected by the \emph{search}, not by the cost model,
whose own separation is only $C_{\mathrm{ins}} = 4.0$ per quarter turn of
claim error.

This is not hypothetical. On one session the 1-move-off decoy is rejected at
beam 4000 and at 16000, and \emph{accepted at beam 64000 at cost 10.17} ---
exactly $6.17 + C_{\mathrm{ins}}$, the constructive bound. Operationally the
system does reject near decoys at the deployed beam, which is what a
verifier has to do; but reporting that as evidence that no cheap wrong story
exists would be the easiest way to oversell this pipeline, and
\cref{eq:decoybound} says one does.

\subsection{The count gate}
\label{sec:count-gate}

\subsubsection{Threat model}

Everything physical reduces to three attacks, and naming them this way
matters because two of them are not swaps, so swap detection cannot see
them.

\begin{center}\small
\begin{tabular}{@{}clp{7.6cm}@{}}
\toprule
& \textbf{Attack} & \textbf{What can catch it} \\
\midrule
A & Substitution --- swap in a solved cube
  & move-count floor; frame-continuity guard; persistent appearance change \\
B & Solving outside the timed window --- follow a solver's output, or stop
    the timer unsolved and solve before scanning
  & move-count floor; the post-stop phantom-rate test \\
C & Video injection / virtual camera
  & \emph{nothing in the video}; challenge--response only \\
\bottomrule
\end{tabular}
\end{center}

\subsubsection{Why counting, and why the metric is the trap}

Move-by-move verification is dead (the cliff above). But that cliff is about
\emph{which} move. Counting discards substitutions entirely and is sensitive
only to insertions and deletions, so it runs on the model's strong axis. MER
is an over-estimate of counting error.

Attack A is defeated trivially: a swapped-in solved cube requires zero moves.
Attack B's first form is the interesting one --- run the scramble through a
solver and execute the answer. Every move genuinely happens; one cube is on
camera throughout; nothing is fake except the thinking.

But a solver's output is \emph{bounded} and a human method is not. So the
floor sits above the longest solution any solver would produce.

\begin{keybox}
\textbf{God's number is 20 in the half-turn metric. This alphabet has no
half-turn symbol}, so \wca{R2} decodes as \wca{R},\wca{R} and everything
here is quarter-turn metric, where God's number is \textbf{26}. A floor set
at ``twenty-something'' would sit \emph{below} the cheat ceiling and catch
nothing.
\end{keybox}

The operative ceiling is not 26 but 30: a cheat runs a two-phase solver, not
an optimal one, and two-phase returns ${\sim}19$--$22$ HTM $\approx 25$--$30$
QTM. The recorded scramble takes corroborate this at 20--28 QTM true.
Constants: solver ceiling 30 QTM, floor on observed moves 32, assumed
efficient-human floor 45 QTM.

\subsubsection{Calibration, and the abstain band}

The gate sees \emph{predicted} moves, and the model under-counts, more so the
faster the solve. Under-counting pushes a cheat further below the floor
(harmless) and a legit solve toward it (the false-DQ risk), so the floor must
be calibrated against worst-case retention, not mean. Retention is measured
by time-warping real sessions to a target turn rate and comparing decoded
count to true count. Taking the \emph{worse} of the two seeds at each level:
at 5\,turns/s retention is 0.88, at 8\,turns/s 0.80, at 10\,turns/s 0.69.

Above some rate the two bands genuinely overlap and no threshold separates
them. That limit is \emph{derived} from the floor rather than hardcoded, so
changing the floor moves it automatically; speed augmentation moved it from
7.11 to 9.62\,turns/s. At 9.62\,turns/s a 45-move solve takes 4.7\,s, against
a ${\sim}3.1$\,s world record --- so the ``too fast to read'' abstain case now
covers only the very fastest humans alive.

\subsubsection{Three tiers, and the order of the checks}

The gate returns \texttt{VERIFIED}, \texttt{REJECTED} or \texttt{REVIEW},
unlike the frame-continuity guard which is deliberately two-tier. Two
conditions here make the evidence \emph{unreadable} rather than
incriminating: too fast to separate, and bad light. Bad light matters
specifically because in low light the \emph{onset detector} loses 40+ points
while the classifier loses ${\sim}5$ --- and the count \emph{is} the
detector's output, so bad light would fabricate a ``too few moves'' rejection
out of a model failure. Abstain conditions are therefore checked
\textbf{before} reject conditions, so a solve is never rejected on evidence
that has just been declared unreadable.

\subsubsection{The hole that remains}

Stated plainly, because it is easy to believe the gate closed more than it
did. \textbf{Make a full solve's worth of plausible moves on camera without
solving, then swap in a solved cube.} That reads as ${\sim}50$ moves, clears
the floor of 32, and the count gate passes it --- eight seconds of flailing
suffices, which would be a world-class time. Nothing about the move count
distinguishes it from a real solve, because the move count is genuinely real;
only the outcome is faked. Frame continuity and appearance own this case
entirely, which is why the count gate can never be the only arm.

Camera injection (attack C) is not addressed by any video analysis. The
defence is protocol: server-issued unique scrambles mean a pre-recorded video
cannot match a fresh challenge, plus challenge--response (demand a specific
face at a random moment). Not built.
